% sections/04-conclusion.tex --- Висновки та перспективи подальшого дослідження.

\citsection{Conclusions and further research}

This paper presented PARCS-Agent, an architectural pattern in which a Model Context Protocol server provides the interface between a language-model agent and a Kubernetes-based PARCS distributed computing cluster. The contribution is positioned at three complementary levels. At the architectural level, a session--layer--worker hierarchy decouples the agent's reasoning steps from cluster orchestration. At the protocol level, four MCP tools --- capability discovery, session compilation, blocking layer execution, and session enumeration --- align with the agent's step-wise decomposition without imposing distributed-computing concepts on the agent's prompt. At the methodological level, runtime compilation of agent-supplied per-worker logic preserves the synchronous coding paradigm at the agent's surface while delivering parallel execution beneath it, thus circumventing the documented difficulty LLMs encounter when generating parallel code directly.

Three design decisions from this work are worth conserving within similar systems. First, the treatment of MCP tool descriptions as load-bearing prompt content reduces the rate of malformed agent calls. Second, the addressing of layer results by identifier rather than by value --- wherein the agent references prior outputs through \texttt{previousLayerId} while the server retrieves the stored payload internally --- bounds the agent's reasoning trace independently of result size and removes a class of state-management concerns from the protocol surface. Third, the distribution of bulk inputs through a content-addressed shared-storage side channel bounds the protocol payload independently of dataset volume and worker count.

Several directions remain open for further investigation. The most immediate is \textit{quantitative evaluation}: a benchmark of compute-bound agent tasks measuring wall-clock duration, solution quality, token consumption, and operational cost relative to a single-threaded sandbox baseline. A companion harness has been developed for this purpose and is reserved for a separate publication. \textit{Multi-language \texttt{IAgentComputation} support} (Python, F\#) would lower the entry barrier for agents whose pre-training distribution emphasises languages other than C\#. A \textit{GPU-backed daemon pool} would extend the architectural pattern to neural-network training and inference workloads. A \textit{cost-aware scheduling} component would enable cloud deployments to balance throughput and operational expenditure without involving the agent in operational concerns. A \textit{formal security model} would be required to address the multi-tenant trust boundary characteristic of any production deployment. The protocol itself appears stable; future development is anticipated to extend the runtime rather than to modify the agent-facing interface.
